\section{An explicit rate for a rational rotation}\label{sec:rational}
Let $\zeta=(3+4i)/5=e^{2\pi i\alpha}$, $0<\alpha<1/4$.
This is not a root of unity: otherwise $\zeta+\zeta^{-1}=6/5$ would
be a rational algebraic integer and hence an integer. Thus $\alpha$
is irrational. No bounded-type assertion about this angle is needed.

\begin{lemma}[Elementary denominator control]\label{lem:rational-denominator}
For every integer $r\ge1$,
\[
 \|r\alpha\|_{\R/\mathbb Z}\ge \frac{5^{-r}}{2\pi}.
\]
Consequently the denominator function of \eqref{eq:denominator-function}
satisfies $Q_\alpha(M)<7\,5^{8M}$ for every integer $M\ge3$.
\end{lemma}
\begin{proof}
The Gaussian integer $(3+4i)^r-5^r$ is nonzero, so its modulus is at
least one. Therefore
\[
 5^{-r}\le |\zeta^r-1|
            \le 2\pi\|r\alpha\|_{\R/\mathbb Z}.
\]
If $q=Q_\alpha(M)$ and $r$ is its preceding convergent denominator,
then $r<8M$ and \eqref{eq:cf} gives
$\|r\alpha\|<1/q$. Combine the two estimates and use $2\pi<7$.
The elementary bound is exponentially weak in $M$, but entirely explicit.
\end{proof}

\begin{theorem}[Rational-angle profile bounds]\label{thm:rational}
For the rotation $R=(1/5)\left(\begin{smallmatrix}3&-4\\4&3\end{smallmatrix}\right)$,
every exact hidden-state profile obeys
\begin{equation}\label{eq:rational-profile}
           \sum_{t<N}2^{-2K_t}5^{-8\,2^{K_t}}\le140.
\end{equation}
In particular,
\begin{equation}\label{eq:rational-peak}
                 N\le140\,2^{2W_N}5^{8\,2^{W_N}}.
\end{equation}
For $N>140$ this implies
\[
 W_N\ge\max\left\{3,
       \log_2\left(\frac19\log_5\frac N{140}\right)\right\}.
\]
In the vector-state class,
$\sum_{t<N}K_t^{-2}5^{-8K_t}\le140$ and, for $N>140$,
$V_N\ge\max\{3,\frac19\log_5(N/140)\}$.
Both classes have an exact rational realization with peak $3(N+1)$.
\end{theorem}
\begin{proof}
Insert Lemma~\ref{lem:rational-denominator} into
\eqref{eq:general-profile}. It gives
\[
 \sum_{t<N}M_t^{-2}5^{-8M_t}
       \le\frac{280\sqrt2}{3}<140,
 \qquad M_t=2^{K_t}.
\]
Each summand is decreasing in $M_t$, proving \eqref{eq:rational-peak}.
For an integer $W\ge1$, $2W\log_5 2\le2^W$; hence
$2\log_5(2^W)+8\,2^W\le9\,2^W$.
Taking logarithms gives the stated hidden lower bound.
For vector states use $M_t=K_t$ and
$2\log_5 V\le V$ for integers $V\ge1$.

For completeness, an exact rational upper machine stores $(s,k)$,
where $s$ is one of the three triangle vertices in
Proposition~\ref{prop:separate} and $k$ is the number of active commands.
Initialize $s$ by the rational barycentric coordinates of $x/10$.
An idle command leaves $(s,k)$ unchanged, and an active command increments
$k$. At epoch $t$ only $0\le k\le t$ is needed. Assign predictive vector
$R^k v_s$ to $(s,k)$ and use its coordinate as the final binary mean.
All assigned vectors have norm less than one, all transitions are
rational, and the row identities hold on every state. The largest
processing alphabet is $3(N+1)$; a held binary answer needs only two
labels in its separate terminal cut.
\end{proof}

The explicit inequalities replace the former quantifier-elimination-only
exclusion for this angle. They do not assert its sharp asymptotic order.
Sharper irrationality estimates for $\alpha$ would improve the denominator
function in the same profile theorem, but none is silently assumed here.
The rational upper machine admits the separately priced fair-bit
implementation of \cite{GTF31}; its atomic label optimum and its sampler
workspace are different quantities.
